\section{Total-product divisibility in the pair-companion algebra}
\label{sec:pair-total-saturation}

This section concerns the precise algebra generated by
$W$ and the even quotient $T'=(T/\pi Z)[1/p]$. Its
source and coefficient maps are defined below. The quotient maps from $H$ to $Q_3$ and $Q_4$ were
constructed in the preceding sections. Here the source
is the displayed free weighted algebra in every degree.

Choose a finite extension $L/\mathbf Q_p$ splitting
the characteristic polynomial of $T'$, and put
$R=O_L[[\gamma-1]][1/p]$. The established rational
module $W_L$ is finite free over $R$. Choose an
$L$-linear section of
\[
 0\longrightarrow N_L\longrightarrow\bar E_L
                  \longrightarrow T'_L\longrightarrow0,
 \qquad N_L=\widehat\bigwedge^2W_L.
\]
Let $A$ be the completed free associative algebra on
$W_L$ in weight one and $T'_L$ in weight two. The
action on $\bar E_L$ defines the Hopf
automorphism
\[
        \gamma t=A_Tt+C_t,\qquad C_t\in N_L,
\]
together with the regular action on $W_L$. All
completed coefficient modules at a fixed degree have
a common bounded denominator in $L$; they are obtained
by taking the indicated compact $O_L$-modules and then
inverting $p$, rather than allowing a separate
denominator for every coordinate.

For a Jordan block with root $\alpha$, length $r$ and
top lift $y$, put $C=(\gamma-\alpha)^ry\in N_L$.
The pair-companion substitution is
\[
 (\gamma-\alpha)^jy\longmapsto
       \frac{C(z_i,z_j)}{(z_iz_j-\alpha)^{r-j}},
                         \qquad0\leq j<r.
\]
The repeated indices on the right denote the two
positions of the token, not its Jordan index.
For a fixed root, the companion residues are
independent over $L$: a relation would give a
$\Lambda_L$-torsion element in $\bar E_L$ with a
nonzero Jordan-socle image in $T'_L$.
Write $q$ for a polynomial annihilating $T'_L$.
In weighted degree $n$ the substitution defines
\[
 \iota_n:A_n\longrightarrow\mathcal B_n^{\oplus c_n},
 \qquad
 \mathcal B_n=O_L[[X_1,\ldots,X_n]][1/p,
                      q(z_iz_j)^{-1}:i<j],\quad z_i=1+X_i,
\]
where the finite array indexes the ordered colors.
It is injective: extract the highest pole of an
initial two-position token, use the independent
companion residues with disjoint tail variables, and
induct on the tail degree; the remaining initial
one-position terms are ordinary coefficient tensors.

\begin{theorem}\label{thm:pair-total-saturation}
For $\beta\in1+\mathfrak m_L$, put
$D_\beta=z_1\cdots z_n-\beta$. For $n=1$ and every
$n\geq3$, and for every $k\geq1$,
\[
 \iota_n(A_n)\cap D_\beta^k\mathcal B_n^{\oplus c_n}
                   =\iota_n((\gamma-\beta)^kA_n).
\]
The same equality holds on the free weighted Lie
subalgebra of $A$.
\end{theorem}

The exclusion of degree two is necessary. At a root
of $q$ its total divisor is already inverted in the
target, whereas the source retains nonzero
coinvariants. The induction below starts with the
separately proved degrees three and four and never
uses that excluded case.

\begin{proof}
Degree one is the usual exact coefficient division.
The initial-token argument below also proves degree
three directly: its total-tail divisor then involves
one variable, is distinct from the pair divisors of
the first-$W$ term, and uses only degree-one source
division. Thus it proves the associative base as
well as the Lie base of
Proposition~\ref{prop:degree-three-saturation}.
We give the additional degree-four argument in
detail, then the exact source induction.

Put $B=A_2$ and $V=W_L\widehat\otimes W_L$. Partition
the words of weight four by their middle cut. There
is an exact $\Gamma$-stable sequence
\begin{equation}\label{eq:pair-four-cut}
 0\longrightarrow B\widehat\otimes_L B
 \longrightarrow A_4
 \longrightarrow W_L\widehat\otimes_L T'_L
                    \widehat\otimes_L W_L\longrightarrow0.
\end{equation}
Its quotient consists of the single token pattern
$W,T',W$. Replacing its middle token by its quadratic
companion lands in the first term, so its induced
action is the diagonal linear action.

Suppose $\iota_4(a)$ vanishes modulo $D_\beta$. On
that hypersurface put $z_4=\beta/(z_1z_2z_3)$.
A middle-pair divisor $z_2z_3-\alpha$ is distinct
from both outer-pair divisor families: their pair
products are $z_1z_2$ and $\beta/(z_1z_2)$.
Consequently a term from $B\widehat\otimes B$ has
no principal part at that middle divisor.

The exact Jordan calculation in the quotient of
\eqref{eq:pair-four-cut} uses $u=z_1z_4$.
For a block of length $r$ with $N=\gamma_T-\alpha$,
the coinvariant relations are
\[
 (u\alpha-\beta)N^jy+uN^{j+1}y=0.
\]
They identify the block cokernel of $uA_T-\beta$
with the singleton coefficient module modulo
$(u\alpha-\beta)^r$, generated by $y$; explicitly
the coefficient vector $F_j(z_1,u)$ maps to
\begin{equation}\label{eq:pair-middle-Jordan}
 \sum_{j=0}^{r-1}F_j(z_1,u)(\beta/u-\alpha)^j
                         \pmod{(\beta/u-\alpha)^r}.
\end{equation}
In the target the middle product is $\beta/u$.
The companion residue depends on the internal
variable $z_2$, whereas the coefficients in
\eqref{eq:pair-middle-Jordan} have the disjoint
singleton variable $z_1$. A finite color and
coefficient minor of the independent residue vectors
detects the largest principal part. Descending through
the pole orders proves that every expression
\eqref{eq:pair-middle-Jordan} vanishes. Division uses
the coordinate $u-\beta/\alpha$ and preserves bounded
coefficients. Thus the quotient image of $a$ is an
$(\gamma-\beta)$-multiple. Lift and subtract
that multiple in $A_4$. We are left with an element
of $B\widehat\otimes B$ having zero total-divisor image.

We compute the required balanced completed tensor
exactly. Let $s=z_1z_2$ and
$\tau_\beta(s)=\beta/s$. The relation
$\gamma_1\gamma_2=\beta$ gives
\begin{equation}\label{eq:pair-balanced-tensor}
 \frac{B\widehat\otimes_L B}
 { (\gamma_1\gamma_2-\beta)(B\widehat\otimes_L B)}
       =B\widehat\otimes_{R,\tau_\beta}B.
\end{equation}
For the generator $\gamma$, this follows by multiplying
the displayed relation by the inverse of the second
$\gamma$. Continuous extension gives balancing for
every bounded series. The involution $\tau_\beta$ is
defined on $R$ because $\beta\in1+\mathfrak m_L$.

The regular pair module $V$ has coordinates
$x=z_1-1$ and $s=z_1z_2$ and is a finite color array of
$O_L[[s-1,x]][1/p]$. Thus
\[
 V\widehat\otimes_{R,\tau_\beta}V
       =O_L[[s-1,x,y]][1/p]^{\oplus c}
\]
with the appropriate finite colors. This is an exact
change of variables; its coefficient map is the
inclusion in the ring with $q(s)$ and $q(\beta/s)$
inverted. In particular its map into
\eqref{eq:pair-balanced-tensor} is injective.
Multiplication by
\[
                           h(s)=q(s)q(\beta/s)
\]
takes \eqref{eq:pair-balanced-tensor} into this regular
submodule. Hence the kernel of its coefficient map
is killed by $h$. It remains to check this kernel
at the finitely many factors of $h$.

Fix $f=s-\alpha$ after splitting these roots in $L$.
Invert the other polynomial factors. The first pair
module has the explicit form
\[
              B=V+\sum_i R f^{-r_i}C_i,
\]
where the sum ranges over the blocks at $\alpha$.
Their residues $C_i\bmod fV$ are independent over $L$.
Choose finitely many internal coefficient and color
coordinates whose matrix on these residues is
invertible. The corresponding matrix $P(s)$ has
determinant $g$ with $g(\alpha)\ne0$. After inverting
this one determinant, the finite-rank coordinate
transformation
\[
 (v_I,v_{\rm rest})\longmapsto
 \bigl(P^{-1}v_I,
               v_{\rm rest}-C_{\rm rest}P^{-1}v_I\bigr)
\]
sends the $C_i$ to distinct coordinate vectors.
Its inverse is the corresponding triangular map.
Both transformations have a fixed common denominator;
their remaining columns are bounded series. The module
$B$ is now a regular coordinate module with finitely
many coordinate lines replaced by $f^{-r_i}R$.

Apply the same transformation to the second factor,
with its twisted action. The balanced completed tensor
is the regular two-internal-variable coordinate module
with selected rows and columns rescaled by
$f^{-r_i-r'_j}$, or by a single one of these factors.
It embeds coordinatewise in the rational coefficient
arrays. This includes the complementary-pair collision:
when the two factors have a pole at the same $f$, their
exponents add. No exclusion of products of eigenvalues
is needed. Every transformation uses finitely many
columns and a fixed denominator, so it acts on the
bounded completion itself. No interchange of arbitrary
products with localization is being presumed.

The kernel therefore vanishes at every factor of $h$.
Since it is killed by $h$, the Bezout identity among
the distinct prime powers makes it zero globally.
This proves injectivity of
\eqref{eq:pair-balanced-tensor} into the specialized
coefficient ring. Returning to
\eqref{eq:pair-four-cut} proves degree-four saturation
for one total divisor. Repetition gives every power.

Now let $n\geq5$. The initial-token decomposition is
the exact $\Gamma$-stable sequence
\begin{equation}\label{eq:pair-initial-token}
 0\longrightarrow P_n=W_L\widehat\otimes A_{n-1}
 \longrightarrow A_n
 \longrightarrow T'_L\widehat\otimes A_{n-2}=Q_n
 \longrightarrow0.
\end{equation}
The right action is diagonal. It has no
$(\gamma-\beta)$-invariants: on each Jordan block
the determinant is a power of
$\alpha\gamma_{A_{n-2}}-\beta$, whose coefficient
action is a nonzero multiplier in a domain.
The adjugate then proves injectivity. Consequently
taking coinvariants in \eqref{eq:pair-initial-token}
is also exact at the left.

There is an exact source evaluation
\begin{equation}\label{eq:pair-source-evaluation}
 (W_L\widehat\otimes A_{n-1})_{\gamma-\beta}
                  \simeq\bigoplus_{\rm colors}A_{n-1},
                 \qquad z_1\longmapsto\beta\gamma_{A_{n-1}}^{-1}.
\end{equation}
Choose a $\Gamma$-stable compact lattice after one
bounded scaling in this fixed degree. The operator
$A=\beta\gamma^{-1}-1$ is topologically nilpotent on
every finite invariant quotient. Evaluation of
$\sum e_jX^j$ at $A$ converges. The identity
\[
 X^j-A^j=(X-A)\sum_{i=0}^{j-1}X^iA^{j-1-i}
\]
gives its exact kernel, with convergent quotient
coefficients. This proves
\eqref{eq:pair-source-evaluation} inside the source,
and bounded scalar extension preserves it.

Suppose $\iota_n(a)$ vanishes on $D_\beta=0$, and put
$s=z_3\cdots z_n$. The pair pole of an initial $T'$
token at $\alpha$ becomes the total-tail divisor
$s-\beta/\alpha$. Since $n-2\geq3$, this divisor is
distinct from every pair divisor in the first-$W$
term and from every proper pair denominator in its
own tail. At its largest pole, the companion residue
is $C(\alpha/z_2,z_2)$, depending only on $z_2$ and
its colors; its coefficient is the tail image on
$s=\beta/\alpha$. The finite constant residue minor
forces every such tail image to vanish. Saturation
in degree $n-2$ therefore puts the source
tail in $(\gamma_A-\beta/\alpha)A_{n-2}$.

The exact Jordan relation in $Q_n$ is
\begin{equation}\label{eq:pair-Jordan-reduction}
 N^jy\otimes(\gamma_A-\beta/\alpha)b
 \equiv-\alpha^{-1}N^{j+1}y\otimes\gamma_A b
                                  \pmod{(\gamma-\beta)Q_n}.
\end{equation}
It follows by expanding
$(\gamma_T\gamma_A-\beta)(N^jy\otimes b)$.
The right side has one lower pole, and is zero at
the end of the Jordan block. Lift each subtraction
to $A_n$; its additional quadratic terms lie in
$P_n$ and have no total-tail pole. Descending through
the finite pole orders and roots shows that the
class of $a$ in $(Q_n)_{\gamma-\beta}$ is zero.

Subtract an $(\gamma-\beta)$-multiple and use
\eqref{eq:pair-source-evaluation}. The remaining
class corresponds to an element of
$\bigoplus A_{n-1}$ whose ordinary coefficient image
in the independent variables $z_2,\ldots,z_n$ is zero.
The already established coefficient injection makes
it zero. This proves saturation at $n$ from saturation
at $n-2$. The bases $3,4$ prove every $n\geq3$,
without using degree two. Cancelling one total
divisor in the coefficient domain and repeating
proves every power $k$.

Finally use the Eulerian primitive projector
\cite[Section~1.1, p.~2]{BandieraSchaetz}
\[
 e=\log^\star(\mathrm{Id})
   =\sum_{j\geq1}\frac{(-1)^{j+1}}j
             (\mathrm{Id}-\eta\epsilon)^{\star j}.
\]
It is a finite sum on every positive weighted degree
of this connected graded cocommutative Hopf algebra,
has image the primitive subspace and is identity
there. Indeed, the coproduct of the identity map is
the convolution product of the maps $a\mapsto a\otimes1$
and $a\mapsto1\otimes a$; these commute by cocommutativity.
Its convolution logarithm therefore satisfies
$\Delta e(a)=e(a)\otimes1+1\otimes e(a)$ for positive
degree $a$. On a primitive element all convolution
powers above the first vanish, so $e(a)=a$.
The finite coefficient formulas extend to the bounded
completion. Naturality makes it commute with the actual
Hopf automorphism $\gamma$, including its quadratic
companion terms. If a primitive $a$ is
$(\gamma-\beta)^k b$ in the associative algebra,
then $a=(\gamma-\beta)^k e(b)$ in the Lie source.
This proves the asserted Lie equality. A projector
depending on ordinary token count would not justify
this step, since that count changes under $\gamma$.
\end{proof}
